\documentclass[10pt,a4paper]{article}
\usepackage[utf8]{inputenc}
\usepackage[T1]{fontenc}
\usepackage[italian]{babel}
\usepackage[default]{sourcesanspro}
\usepackage{inconsolata}
\usepackage{amsmath,microtype,booktabs,tabularx,array,enumitem}
\usepackage[table]{xcolor}
\usepackage{geometry,fancyhdr,lastpage,tcolorbox,listings,hyperref}
\tcbuselibrary{skins}
\geometry{left=19mm,right=19mm,top=17mm,bottom=16mm,
  headheight=21pt,headsep=5mm,footskip=9mm,includeheadfoot}
\definecolor{Navy}{HTML}{16334A}
\definecolor{Teal}{HTML}{147D83}
\definecolor{Ink}{HTML}{253947}
\definecolor{Muted}{HTML}{60717D}
\definecolor{Mist}{HTML}{EEF4F7}
\definecolor{Rule}{HTML}{D5E1E7}
\definecolor{Warm}{HTML}{A35E28}
\definecolor{Sand}{HTML}{FAF3E9}
\color{Ink}
\setlength{\parindent}{0pt}
\setlength{\parskip}{5pt}
\setlength{\emergencystretch}{2em}
\linespread{1.08}
\setlist[itemize]{leftmargin=5mm,itemsep=3pt,topsep=4pt,parsep=0pt}
\setlist[enumerate]{leftmargin=6mm,itemsep=3pt,topsep=4pt,parsep=0pt}
\renewcommand{\arraystretch}{1.18}
\setlength{\tabcolsep}{5pt}
\hypersetup{colorlinks=true,linkcolor=Teal,urlcolor=Teal,
  pdftitle={Decentralised Lending Service - Capire e spiegare il codice},
  pdfauthor={Marco Pernisco},pdfsubject={Guida individuale per l'esame orale}}
\pagestyle{fancy}\fancyhf{}
\fancyhead[L]{\fontsize{8}{10}\selectfont\color{Muted}DECENTRALISED LENDING SERVICE}
\fancyhead[R]{\fontsize{8}{10}\selectfont\color{Muted}CAPIRE E SPIEGARE IL CODICE}
\fancyfoot[L]{\fontsize{8}{10}\selectfont\color{Muted}Marco Pernisco / 683674}
\fancyfoot[R]{\fontsize{8}{10}\selectfont\color{Muted}\thepage\ /\ \pageref*{LastPage}}
\renewcommand{\headrulewidth}{0.3pt}
\renewcommand{\headrule}{\hbox to\headwidth{\color{Rule}\leaders\hrule height \headrulewidth\hfill}}
\newcommand{\code}[1]{\textcolor{Navy}{\texttt{\small #1}}}
\newcommand{\lead}[1]{\par\vspace{2mm}{\large\bfseries\color{Navy}#1\par}\vspace{0.5mm}}
\newcounter{guidechapter}
\newcommand{\chapterpage}[3]{\clearpage
  \stepcounter{guidechapter}\pdfbookmark[0]{#2}{chapter\theguidechapter}
  {\small\bfseries\color{Teal}#1\par}\vspace{1mm}
  {\fontsize{24}{28}\selectfont\bfseries\color{Navy}#2\par}\vspace{2mm}
  {\color{Muted}#3\par}\vspace{3mm}}
\tcbset{panel/.style={enhanced,sharp corners,boxrule=0.4pt,colframe=Rule,
  colback=Mist,left=3mm,right=3mm,top=2mm,bottom=2mm,before skip=5pt,after skip=5pt}}
\newtcolorbox{keybox}[1]{panel,borderline west={1.5pt}{0pt}{Teal},
  title={#1},fonttitle=\bfseries,fontupper=\raggedright,coltitle=Navy,colbacktitle=Mist,titlerule=0pt}
\newtcolorbox{warningbox}[1]{panel,colback=Sand,colbacktitle=Sand,
  borderline west={1.5pt}{0pt}{Warm},title={#1},fonttitle=\bfseries,fontupper=\raggedright,coltitle=Warm,titlerule=0pt}
\newcommand{\oral}[2]{\begin{keybox}{Il professore potrebbe chiedere: #1}#2\end{keybox}}
\lstdefinelanguage{Solidity}{morekeywords={address,bool,bytes,calldata,constant,
  constructor,contract,else,emit,error,event,external,false,for,function,if,
  immutable,import,interface,internal,is,library,mapping,memory,modifier,new,
  onlyOwner,override,payable,private,public,pure,receive,return,returns,revert,
  storage,string,struct,true,uint160,uint256,view,virtual},sensitive=true,
  morecomment=[l]{//},morestring=[b]"}
\lstset{basicstyle=\ttfamily\fontsize{8.5}{10.5}\selectfont,
  keywordstyle=\color{Teal}\bfseries,commentstyle=\color{Muted},
  stringstyle=\color{Warm},numbers=left,numberstyle=\tiny\color{Muted},
  numbersep=7pt,xleftmargin=8mm,framexleftmargin=6mm,
  backgroundcolor=\color{Mist},frame=single,rulecolor=\color{Rule},
  framerule=0.3pt,framesep=2mm,columns=fullflexible,keepspaces=true,
  breaklines=true,breakatwhitespace=false,showstringspaces=false,
  aboveskip=4pt,belowskip=6pt,tabsize=4}
\newcommand{\sourcecode}[4]{\par
  {\footnotesize\color{Muted}\path{#1}\quad /\quad righe #2--#3\par}
  \lstinputlisting[language=#4,firstline=#2,lastline=#3,firstnumber=#2]{../#1}}
\newcommand{\qa}[2]{\par\vspace{1mm}\textbf{\textcolor{Navy}{#1}}\par #2\par}

\begin{document}
\pdfbookmark[0]{Come usare la guida}{inizio}
{\small\bfseries\color{Teal}GUIDA INDIVIDUALE / PREPARAZIONE ALL'ORALE}\par
\vspace{5mm}
{\fontsize{31}{35}\selectfont\bfseries\color{Navy}Il progetto,\\spiegato con il codice\par}
\vspace{4mm}
{\Large Una lettura guidata per capire che cosa succede,\par
perché succede e come raccontarlo al professore.\par}
\vspace{3mm}
Marco Pernisco\quad\textcolor{Muted}{683674 / P2P Systems and Blockchains Lab}

\lead{Da dove partiamo}
Immagina una cassa comune. Alcune persone vi depositano ETH; un'altra persona
chiede un prestito. Chi ha messo i fondi vota se concederlo. Se il prestito viene
approvato, il denaro esce dalla cassa e arriva al richiedente. Quando rientra,
una parte restituisce il capitale e un'altra paga gli interessi.

Il nostro progetto traduce questa storia in contratti Solidity e programmi
Python. La parte difficile non è ricordare i nomi: è capire come ogni istruzione
cambia i soldi disponibili, i crediti e lo stato di un prestito.

\begin{keybox}{Come studiare ogni estratto}
Leggi prima la situazione descritta. Poi guarda il codice e chiediti:
\textbf{chi lo chiama, che cosa controlla, quali dati cambia e a chi invia ETH?}
Infine prova a rispondere alla domanda nel riquadro senza leggere la soluzione.
La risposta deve spiegare il motivo del codice, non tradurre soltanto le parole.
\end{keybox}

\lead{Tre percorsi di lettura}
\begin{tabularx}{\linewidth}{@{}>{\bfseries}p{24mm}X@{}}
\toprule
Il prestito & Persone e saldi; depositi; voto; erogazione; rate; default e compensazione.\\
I programmi & Oracle Bitcoin, saldi UTXO, attesa iniziale, voter e gestione degli eventi.\\
La discussione & Setup, proxy, sicurezza, test, gas, strategia economica e demo.\\
\bottomrule
\end{tabularx}

\begin{warningbox}{Codice vero, non pseudocodice}
Gli estratti sono letti direttamente dai file del progetto durante la compilazione
del PDF. I numeri a sinistra sono le righe dei sorgenti. Mostriamo porzioni di
funzioni: il contesto e gli altri controlli sono spiegati nel testo. Se i file
cambiano, i riferimenti vanno ricontrollati. La guida è separata dal report di
consegna e non modifica le regole dei contratti.
\end{warningbox}

\chapterpage{01 / LA STORIA COMPLETA}{Chi mette i fondi e chi li riceve}{Prima distinguiamo le persone dai programmi che usano.}
Chiamiamo \textbf{Alice} la persona che deposita nel pool e \textbf{Bob} quella
che chiede il prestito. Questi nomi servono solo agli esempi; il codice identifica
le persone con indirizzi Ethereum.

Alice è un \emph{contributor}: mette ETH a disposizione, vota e può ritirare la
parte non impegnata in prestiti. Bob è l'\emph{applicant}, cioè il richiedente;
quando riceve il prestito lo chiamiamo anche \emph{borrower}. Una stessa persona
può ricoprire entrambi i ruoli. Esistono poi l'owner, che amministra il pool,
e l'operatore oracle, che pubblica il dato Bitcoin.

\lead{La sequenza, raccontata senza formule}
Bob indica un indirizzo Bitcoin e paga una richiesta di aggiornamento. Un
programma Python ricostruisce quanto quell'indirizzo possiede nello snapshot
storico richiesto e scrive il risultato su Ethereum. Bob propone il prestito;
Alice vota. Passati i blocchi previsti, Bob chiede al pool di risolvere la proposta.

Se voto, disponibilità ETH e controllo BTC permettono il prestito, il pool crea
un contratto dedicato e gli passa il capitale. Il contratto lo invia a Bob.
Bob paga le rate a quel contratto, che comunica al pool come dividerle. Se Bob
non paga in tempo, il prestito può essere dichiarato fallito. Alice può chiedere
una compensazione, ma solo nei limiti della riserva realmente esistente.

\lead{Dove cercare i vari passaggi}
\begin{tabularx}{\linewidth}{@{}p{46mm}X@{}}
\toprule File & In parole semplici\\\midrule
\code{FundingPool.sol} & La cassa comune e i suoi registri: depositi, voti e riparti.\\
\code{LoanContract.sol} & La scheda di un singolo prestito, con termini e scadenza.\\
\code{BitcoinOracle.sol} & Il registro Ethereum dei saldi BTC comunicati dall'operatore.\\
\code{FundingPoolProxy.sol} & L'indirizzo stabile del pool, anche quando cambia la logica.\\
\code{oracle\_service.py} & Legge Bitcoin e risponde alle richieste di saldo.\\
\code{auto\_voter.py} & Vota sempre YES con l'account configurato.\\
\code{simulate.py} & Fa compiere agli account la sequenza della dimostrazione.\\
\bottomrule
\end{tabularx}

\oral{perché non fa tutto un unico contratto?}{Il pool gestisce i fondi comuni;
ogni contratto di prestito gestisce una singola operazione. Il dato Bitcoin
richiede invece un servizio esterno, perché il contratto Ethereum non apre da
solo i file Bitcoin del computer. Separare questi compiti rende più chiaro chi
può modificare ciascun dato.}

\chapterpage{02 / LEGGERE SOLIDITY}{Le parole che trovi nel codice}{Pochi concetti di sintassi per non bloccarti davanti a una funzione.}
\sourcecode{contracts/FundingPool.sol}{70}{75}{Solidity}
Un \code{mapping} è un registro indicizzato da una chiave. Nella riga 70 la
chiave è un indirizzo: per ogni persona salvo un importo. Nella riga 74 ci sono
due chiavi: prima l'identificativo della proposta, poi l'indirizzo del votante.
Così posso chiedere: ``Alice ha già votato sulla proposta numero 2?''.

\code{address} rappresenta un indirizzo Ethereum; \code{uint256} è un intero
senza segno; \code{bool} vale vero o falso. Gli importi ETH sono interi in
\textbf{wei}: 1 ETH corrisponde a $10^{18}$ wei. Il contratto non usa numeri con
la virgola per i pagamenti.

\sourcecode{contracts/FundingPool.sol}{141}{148}{Solidity}
Qui \code{external} permette di chiamare la funzione dall'esterno.
\code{payable} permette di accompagnare la chiamata con ETH.
\code{msg.sender} è chi sta chiamando il contratto, mentre \code{msg.value} è
quanto ETH accompagna questa chiamata, espresso in wei.

\code{notTerminated} è un \emph{modifier}: esegue un controllo aggiuntivo prima
del corpo della funzione. \code{revert} interrompe l'operazione e annulla gli
effetti della transazione; il gas già consumato non viene magicamente azzerato.
\code{emit} produce un evento nei log: è una notifica che Python può leggere,
non una chiamata diretta al programma Python. Nel corpo di un modifier,
\code{\_} segna il punto in cui viene eseguita la funzione protetta.

\lead{Altre parole che incontreremo}
\code{storage} indica dati persistenti del contratto: un riferimento storage
permette di modificarli direttamente. \code{memory} indica dati temporanei della
chiamata. \code{view} indica che la funzione non modifica lo stato; \code{pure}
che non legge né modifica lo stato del contratto. \code{immutable} fissa un dato
alla costruzione del contratto. \code{internal} serve alla logica interna o
ereditata, non è un normale ingresso pubblico per l'utente.

\oral{public significa che tutti possono modificare il saldo?}{No.
Per questo mapping il compilatore genera una funzione di lettura. Le modifiche
devono passare per il codice delle funzioni che le permettono. Visibilità del
dato e autorizzazione a cambiarlo sono due cose diverse.}

\chapterpage{03 / I TRE SALDI}{Depositato non vuol dire disponibile}{Il punto più importante della contabilità del progetto.}
Alice deposita 10 ETH. Finché non presta nulla, può ritirare 10. Se 4 vengono
erogati a Bob, Alice ha ancora una posizione complessiva di 10, ma può ritirarne
solo 6: gli altri 4 sono impegnati nel prestito.

\sourcecode{contracts/FundingPool.sol}{172}{174}{Solidity}
La funzione legge il totale registrato per Alice e sottrae il capitale bloccato.
Non interroga il saldo del wallet di Alice e non trasferisce ETH. Calcola la
parte della sua posizione che il pool può ancora utilizzare o restituirle.

\begin{tabularx}{\linewidth}{@{}Xrrr@{}}
\toprule Situazione (ETH) & Totale & Bloccato & Disponibile\\\midrule
Alice deposita & 10 & 0 & 10\\
Bob riceve un prestito di 4 & 10 & 4 & 6\\
Bob restituisce 2 di capitale & 10 & 2 & 8\\
Bob restituisce gli ultimi 2 & 10 & 0 & 10\\\bottomrule
\end{tabularx}

\lead{Perché il totale non sale quando Bob rimborsa?}
Perché quei 4 ETH erano già parte della posizione di Alice. Il rimborso cambia
la loro condizione: da impegnati tornano disponibili. Se aggiungessi anche 4 al
totale, dopo il rimborso Alice risulterebbe creditrice di 14 ETH pur avendone
depositati 10, senza nemmeno contare gli interessi. Sarebbe un doppio conteggio.

Nel pool esistono anche due somme separate: il \textbf{gain}, cioè gli interessi
da incassare, e la \textbf{riserva di compensazione}, usata per le perdite.
Nessuna delle due aumenta automaticamente il capitale disponibile per nuovi
prestiti. Il saldo ETH del contratto e il capitale totale registrato per le
persone, quindi, non sono la stessa misura.

\begin{keybox}{Una regola da ricordare con un esempio, non solo come formula}
Totale = disponibile + bloccato. Nel nostro esempio: $10=6+4$.
Nei normali flussi, gli ETH effettivamente nel pool coprono capitale disponibile,
riserva e gain non incassati. Il capitale bloccato è già stato erogato: non è
ancora denaro liquido dentro il contratto.
\end{keybox}

\oral{perché non posso ritirare anche il bloccato?}{Perché il pool lo ha già
prestato. La variabile rappresenta quanto resta da recuperare, non ETH ancora
disponibili in cassa. Per ritirarlo deve prima rientrare il capitale, oppure
deve avvenire un recupero secondo le regole della compensazione.}

\chapterpage{04 / DEPOSITI E PRELIEVI}{Come entrano ed escono gli ETH}{Seguire l'ordine: controlli, aggiornamenti e trasferimento.}
\sourcecode{contracts/FundingPool.sol}{141}{148}{Solidity}
Il deposito usa l'importo della chiamata, non un parametro dichiarato a parole.
La riga 142 rifiuta un importo minore o \emph{uguale} al minimo: servono
\textbf{più di 100.000 wei}. Poi aggiorna il totale della persona e la disponibilità
del pool. La registrazione aggiunge la persona all'elenco dei contributor solo
se non c'è già; l'evento permette ai client di osservare il deposito.

\sourcecode{contracts/FundingPool.sol}{153}{166}{Solidity}
Il prelievo prima verifica che l'importo sia positivo e non superi il disponibile.
Poi riduce i registri, e soltanto dopo invia gli ETH. Con 10 totali e 4 bloccati,
Alice può prelevare al massimo 6. Se preleva 2, rimangono 8 totali, 4 bloccati e
4 disponibili.

Nella riga 163 \code{payable(msg.sender)} indica il destinatario; la
\code{call} invia \code{amount} wei. Il booleano \code{sent} dice se la chiamata
è riuscita. Se fallisce, il \code{revert} annulla anche le modifiche ai registri:
Alice non perde il credito senza ricevere il denaro.

\oral{perché aggiorni i saldi prima della call?}{Perché il destinatario potrebbe
essere un contratto ed eseguire altro codice appena riceve ETH. Aggiornando
prima, una chiamata ricorsiva non vede più lo stesso saldo disponibile da
prelevare. È la difesa principale mostrata nel nostro test di reentrancy.}

\chapterpage{05 / CREARE E VOTARE}{La proposta non è ancora il prestito}{Salvare i termini e registrare una scelta sono operazioni diverse dall'erogazione.}
Bob deve chiedere un prestito a proprio nome, indicare un indirizzo BTC non
vuoto, un capitale positivo, un tasso tra 1 e 100 e una durata positiva.
Nel codice il parametro del voto deve essere \textbf{esattamente 12}: la
risoluzione è possibile dopo che sono trascorsi quei 12 blocchi.

\sourcecode{contracts/FundingPool.sol}{460}{474}{Solidity}
La riga 467 assegna un nuovo numero alla proposta. La riga 468 prende un
riferimento \code{storage}: le assegnazioni successive modificano la proposta
salvata nel contratto, non una copia temporanea. La scadenza del voto è espressa
come altezza di blocco, non come data o secondi dell'orologio.

\sourcecode{contracts/FundingPool.sol}{490}{502}{Solidity}
Si rifiutano proposte inesistenti, già risolte e voti duplicati. Il controllo
alla riga 498 usa il \emph{totale}, quindi un contributor con tutto il capitale
bloccato può votare. Il suo peso corrente può però essere zero. La coppia di
mapping salva separatamente ``ha votato'' e ``ha scelto YES o NO''.

\oral{perché servono due mapping e non basta voteSupport?}{Un booleano non
impostato vale già false. Senza \code{hasVoted} non distinguerei chi ha scelto
NO da chi non ha ancora votato, né potrei impedire correttamente un secondo
voto. Per decidere l'esito entrambi pesano NO, ma sono situazioni diverse.}

\chapterpage{06 / DECIDERE L'ESITO}{Il voto viene ripesato alla fine}{La scelta resta la stessa, ma il capitale disponibile può cambiare.}
Immagina Alice con 6 ETH disponibili e Carlo con 4. Alice vota YES e Carlo
non vota: in quel momento Alice avrebbe la maggioranza. Prima della risoluzione,
però, Alice ritira 4 ETH. Le disponibilità diventano 2 e 4: il suo YES non ha
più la maggioranza. Il codice deve guardare la situazione corrente.

\sourcecode{contracts/FundingPool.sol}{524}{538}{Solidity}
Il ciclo somma soltanto la disponibilità corrente di chi ha votato YES.
La riga 533 considera NO tutto il resto: sia i contrari, sia chi non ha votato.
Il controllo alla riga 535 richiede più di metà del capitale disponibile.
Una parità non basta. Anche quando la somma è dispari, il confronto intero
funziona: su 5 wei totali, 3 YES bastano e 2 no.

\lead{Perché non uso solo i pesi registrati al momento del voto?}
\code{voteOnProposal()} aggiorna anche i conteggi della proposta, ma quei valori
non sono il risultato definitivo: qui vengono ricalcolati e sovrascritti.
Nel frattempo possono esserci stati depositi, prelievi e rimborsi di altri
prestiti. La regola del progetto richiede i saldi disponibili alla risoluzione.

Nel codice, \code{votingDeadlineBlock} indica quando Bob può risolvere:
dal blocco di creazione più 12 in poi. Non chiude il voto da solo: si può votare
finché la proposta non viene risolta. Solo Bob può risolverla.
Una risoluzione anticipata fallisce; il rifiuto economico invece chiude la
proposta, emette un evento e restituisce l'indirizzo zero.

\oral{a cosa serve executed = true anche quando rifiuti?}{A rendere definitiva
quella risoluzione. Il nome del campo non significa necessariamente ``denaro
erogato'': significa proposta già trattata. Senza chiuderla, Bob potrebbe
provare ripetutamente a farla passare dopo aver cambiato la situazione.}

\chapterpage{07 / TROVARE IL CAPITALE}{Da quali persone arriva il prestito}{Prima verifico la capacità; poi blocco le quote proporzionali.}
\sourcecode{contracts/FundingPool.sol}{547}{553}{Solidity}
L'oracle restituisce satoshi. Il pool li converte in wei con il cambio fisso
di \textbf{30 ETH per BTC}. \code{1 ether} è una costante di unità di Solidity:
vale $10^{18}$ wei, non una richiesta di inviare un ETH.
Con 0,04 BTC la capacità è 1,2 ETH: una richiesta di 1 può passare questo controllo,
una di 2 no. Devono comunque essere sufficienti anche i fondi ETH del pool.

\sourcecode{contracts/FundingPool.sol}{262}{273}{Solidity}
Queste righe sono dentro il ciclo sui contributor. Se il pool ha 60 ETH
disponibili, divisi in 10, 20 e 30, per un prestito di 12 le quote sono 2, 4 e 6.
\code{Math.mulDiv} calcola il prodotto diviso per il totale e, in questa forma,
arrotonda per difetto. Tutti i contributor disponibili partecipano, non solo
chi ha votato YES.

La quota aumenta il bloccato della persona e diminuisce la disponibilità comune.
Vengono salvate due copie con significati diversi: \code{loanContributions}
è il credito che si ridurrà con i recuperi; \code{initialLoanContributions}
conserva la quota iniziale per ordine di rimborso e interessi.

\oral{perché sommi totalLocked invece di erogare la richiesta?}{Perché le quote
intere possono perdere un resto. Con disponibilità 100.001 e 100.002 wei e
richiesta di 3, ogni quota è 1: erogo 2, non 3. I controlli di capacità vengono
prima effettuati sulla richiesta; il contratto del prestito usa poi il capitale
effettivo, che può essere inferiore per arrotondamento.}

\chapterpage{08 / IL CONTRATTO DEL PRESTITO}{Una scheda separata per Bob}{Il pool crea il contratto e gli dà i fondi per l'erogazione.}
\sourcecode{contracts/FundingPool.sol}{563}{575}{Solidity}
\code{new LoanContract} distribuisce una nuova istanza su Ethereum.
Le passo indirizzo del pool, borrower, capitale effettivo, tasso, durata e
percentuale di collateral corrente. Ogni prestito ha quindi i propri termini.
Il pool registra l'indirizzo come autorizzato prima di finanziarlo.

\sourcecode{contracts/LoanContract.sol}{91}{98}{Solidity}
\code{onlyPool} permette il finanziamento solo al pool associato. Il prestito
deve essere nello stato Created e ricevere esattamente il capitale previsto.
Passa a Funded e invia gli ETH a Bob. Se l'invio fallisce, la transazione
si annulla: non resta un finanziamento registrato senza l'erogazione riuscita.

\lead{Qui msg.sender cambia lungo la catena delle chiamate}
Bob chiama il pool per risolvere: nel pool il chiamante è Bob. Il pool chiama
\code{fundLoan()}: dentro il prestito il chiamante è il pool, non Bob.
La variabile \code{borrower} serve a ricordare chi deve ricevere gli ETH e
chi potrà rimborsare. \code{msg.sender} non indica sempre l'utente iniziale.

Il costruttore fissa \code{dueBlock = block.number + durationBlocks}.
Fissa anche il collateral del prestito: se nasce al 50\%, continuerà a usare
il 50\% anche quando il parametro globale del pool cambierà.
Gli stati principali sono Created, Funded, Repaid e Defaulted: proposta e
prestito sono due oggetti diversi, con momenti diversi del ciclo di vita.

\oral{perché controlli authorizedLoanContracts?}{Per riconoscere i contratti
che il pool ha davvero creato. Più avanti il prestito chiamerà funzioni che
sbloccano fondi e registrano default: non posso permettere che un indirizzo
qualsiasi si presenti come un prestito e alteri quei conti.}

\chapterpage{09 / CAPIRE UNA RATA}{Che cosa paga davvero Bob}{Prima facciamo i conti in ETH; poi leggiamo il calcolo Solidity.}
Bob riceve 1 ETH al 5\%. Deve restituire 1 ETH di capitale e 0,05 di interesse,
per un totale di 1,05. Il tasso è sul prestito intero: non stiamo applicando
un tasso annuale né un interesse composto.

\sourcecode{contracts/LoanContract.sol}{79}{86}{Solidity}
\code{dueAmount()} restituisce il debito totale previsto, non il residuo dopo
le rate già pagate. L'interesse è una divisione intera; il capitale effettivamente
erogato è il punto di partenza del calcolo.

Se Bob paga due rate da 0,525 ETH e il collateral del prestito è 50\%, ciascuna
rata contiene 0,5 di capitale e 0,025 di interesse. Metà dell'interesse va alla
riserva e metà è il gain del contributor. Il collateral \textbf{non} trattiene
il 50\% del capitale prestato.

\begin{tabularx}{\linewidth}{@{}Xrr@{}}
\toprule Destinazione & Una rata & Totale due rate\\\midrule
Capitale restituito & 0,5000 & 1,0000\\
Riserva & 0,0125 & 0,0250\\
Gain & 0,0125 & 0,0250\\\midrule
Pagamento di Bob & 0,5250 & 1,0500\\\bottomrule
\end{tabularx}

\lead{Quanto della rata accetto come debito?}
\sourcecode{contracts/LoanContract.sol}{117}{121}{Solidity}
\code{paidBefore} somma capitale e interessi già contabilizzati.
\code{debt - paidBefore} è il debito residuo. Se Bob deve ancora 0,525 ma manda
0,6, accetto 0,525 come rimborso del debito; gli altri 0,075 sono un'eccedenza
e finiranno nella riserva, non vengono restituiti a Bob.

\oral{perché non metti tutti i soldi della rata nel disponibile?}{Perché solo
il capitale ricostituisce i fondi prestabili. Gli interessi hanno due destinazioni
separate: gain dei contributor e riserva. Sommarli al disponibile falserebbe
la contabilità e permetterebbe di prestare somme già dovute ad altro titolo.}

\chapterpage{10 / IL CALCOLO DELLE RATE}{Perché guardo quanto è stato pagato}{Non riparto da zero a ogni rata: aggiorno un totale progressivo.}
Il problema è l'arrotondamento: dividere separatamente ogni pagamento può
produrre risultati diversi a seconda di quante rate Bob sceglie. Per questo
calcoliamo prima quale dovrebbe essere il totale raggiunto dopo l'ultimo
pagamento, poi sottraiamo quello che avevamo già registrato.

\sourcecode{contracts/LoanContract.sol}{125}{135}{Solidity}
Nella riga 125 \code{paidAfter} è quanto Bob ha pagato complessivamente, limitato
al debito. La proporzione fra capitale e debito ci dice quanto di quel totale
è capitale. La riga 126 assegna il resto all'interesse. Le righe 127--128
calcolano solo la parte nuova, che appartiene alla rata attuale.

Con il prestito di 1 ETH, dopo la prima rata abbiamo 0,5 di capitale e 0,025 di
interesse. Dopo la seconda, i totali diventano 1 e 0,05. La seconda rata
contabilizza le differenze: $1-0{,}5=0{,}5$ e $0{,}05-0{,}025=0{,}025$.
Non deve distribuire di nuovo i totali già raggiunti con la prima rata.

\lead{La stessa idea viene usata per il collateral}
Le righe 130--131 calcolano quanto dovrebbe esserci complessivamente nella
riserva per gli interessi pagati, e tolgono la parte già conteggiata prima.
La riga 132 assegna il resto dell'interesse al gain. Solo dopo aggiorniamo
\code{repaidPrincipal} e \code{repaidInterest}: i vecchi valori sono serviti
per calcolare le differenze. L'eccedenza viene aggiunta alla quota di riserva.

La formula matematica è $B(C)=\lfloor CP/(P+I)\rfloor$: $C$ è il pagato totale,
$P$ il capitale e $I$ l'interesse previsto. Non occorre recitarla senza capirla:
è una proporzione applicata al totale raggiunto. È una scelta implementativa
documentata per rendere coerenti rate e arrotondamenti.

\oral{e se il professore ti indica principalAfter - repaidPrincipal?}{Direi:
``Il primo valore è il capitale che risulta rimborsato dopo questa rata; il
secondo è quello già rimborsato prima. La differenza è il nuovo capitale da
sbloccare adesso. Se distribuissi principalAfter per intero, conterei due
volte le rate precedenti.''}

\chapterpage{11 / COLLEGARE I CONTRATTI}{Come il rimborso torna al pool}{Prima arriva il denaro; poi vengono aggiornati i registri comuni.}
Bob chiama \code{repay()} sul contratto del proprio prestito. Dopo aver
calcolato capitale, gain e quota di riserva, il prestito deve trasferire gli
ETH al pool e spiegargli come dividerli. Non basta aggiornare i contatori
del solo prestito: i saldi dei contributor si trovano nel pool.

\sourcecode{contracts/LoanContract.sol}{137}{145}{Solidity}
La prima condizione chiude con successo soltanto un prestito ancora Funded
che abbia ricevuto tutto il capitale e l'interesse. Se era già Defaulted,
quel controllo non lo trasforma in Repaid: è il motivo concreto per cui
un pagamento tardivo non cancella il fallimento.

La riga 141 invia al pool tutto il denaro della rata, compresa un'eventuale
eccedenza. Il \code{receive()} del pool accetta l'invio soltanto da prestiti
autorizzati. Poi l'interfaccia \code{IFundingPool} permette di chiamare
\code{handleRepayment()}: non crea un altro pool, descrive come chiamare quello
all'indirizzo già noto. In questo callback il chiamante è il prestito.

\sourcecode{contracts/FundingPool.sol}{398}{410}{Solidity}
Il controllo iniziale protegge questi aggiornamenti: un estraneo non può
chiamarli inventando capitale da sbloccare. Poi il pool usa tre strade:
restituisce il capitale, accredita il gain e aumenta la riserva.
I parametri descrivono gli ETH già ricevuti, non inviano tre nuovi pagamenti.

\oral{perché ci sono sia una call con ETH sia handleRepayment?}{La prima sposta
il denaro dal prestito al pool; la seconda aggiorna la sua contabilità secondo
la divisione calcolata. Sono due ruoli distinti dentro la stessa transazione.
Se il callback fallisce, il revert annulla anche il trasferimento precedente
e gli altri effetti di quella transazione.}

\chapterpage{12 / RESTITUIRE IL CAPITALE}{Chi recupera per primo}{Il capitale segue una priorità; il gain segue una proporzione.}
Se Alice ha finanziato 6 ETH e Carlo 4, il capitale torna prima ad Alice.
Quando rientrano 2 ETH di \emph{capitale}, il credito di Alice scende da 6 a 4,
quello di Carlo resta 4. La priorità è data dalla quota iniziale, non viene
ricalcolata cambiando ordine dopo ogni rata.

\sourcecode{contracts/FundingPool.sol}{207}{224}{Solidity}
Il ciclo usa un elenco già ordinato. \code{remaining} è il capitale ancora da
assegnare; \code{lockedForLoan} è il credito residuo di questa persona su questo
prestito. Il valore \code{refund} prende il minore dei due. Se restano 3 da
assegnare e Alice ha credito 1, ad Alice va 1 e il ciclo continua con gli altri 2.

Le righe 214--215 riducono sia il credito su quel prestito sia il bloccato totale
della persona. La riga 217 aumenta il disponibile del pool. Non aumentiamo
\code{totalContributed}: il capitale era già contato. Un eventuale capitale
non più assegnabile, ad esempio perché un credito è già stato compensato,
finisce nella riserva.

\sourcecode{contracts/ContributorOrder.sol}{30}{33}{Solidity}
Questa funzione confronta due posizioni: prima viene la quota iniziale maggiore;
in parità, l'indirizzo numericamente maggiore. \code{uint160} permette il
confronto numerico degli indirizzi. Il verso decrescente è la scelta documentata
per rendere deterministico lo spareggio.

\oral{a cosa serve loanOrderPrepared?}{A non rifare l'ordinamento a ogni rata.
Un merge sort prepara la priorità una volta per prestito; il flag ricorda che
il lavoro è già fatto. La preparazione ritardata serve anche ai prestiti
preesistenti che non avevano ancora quel flag impostato.}

\chapterpage{13 / GLI INTERESSI}{Accreditare non è inviare denaro}{Il gain diventa un credito separato, poi Alice lo incassa.}
Con quote iniziali 6 e 4, Alice ha diritto al 60\% del gain e Carlo al 40\%.
Se il gain da dividere è 0,025 ETH, riceveranno crediti di 0,015 e 0,010.
Anche se Alice ha recuperato prima il capitale, la sua quota iniziale rimane
quella usata per gli interessi.

\sourcecode{contracts/FundingPool.sol}{304}{313}{Solidity}
Queste righe sono dentro il ciclo sui finanziatori. Ogni quota viene aggiunta
a \code{pendingInterestGain}: è un aggiornamento del registro, non un bonifico.
\code{disbursed} somma ciò che è stato accreditato, nonostante il nome possa
far pensare a denaro già inviato. Le divisioni intere possono lasciare qualche
wei: il resto viene aggiunto alla riserva.

\sourcecode{contracts/FundingPool.sol}{330}{338}{Solidity}
Qui invece avviene l'invio. Si legge il credito del chiamante, lo si azzera
prima della chiamata esterna e si trasferisce l'importo al suo wallet.
Se il trasferimento fallisce, il revert ripristina anche il credito.

\lead{Perché separare i due momenti?}
Durante il rimborso non inviamo ETH a tutti i contributor uno per uno. Ognuno
ritira il proprio gain con una chiamata esplicita: è il modello \emph{pull},
cioè ``vengo a prendere ciò che mi spetta''. Si separa il guadagno dal capitale
prestabile e non si delega a ogni destinatario l'esito della distribuzione
durante una rata.

\oral{perché azzeri pendingInterestGain prima di inviare ETH?}{Per impedire che
un contratto destinatario rientri e incassi di nuovo lo stesso credito. Se
richiama la funzione, il credito è già zero. Se invece l'invio fallisce,
l'intera operazione viene annullata e il credito non è perso.}

\chapterpage{14 / IL FALLIMENTO}{Una scadenza non esegue codice}{Il default viene registrato quando qualcuno manda una transazione valida.}
Supponiamo che il prestito scada al blocco 100. Un rimborso completo nel blocco
100 è ancora puntuale. Dal 101 è tardivo. Il nodo non chiama automaticamente
una funzione allo scoccare del blocco: finché non avviene una transazione,
lo stato salvato può ancora dire Funded.

\sourcecode{contracts/LoanContract.sol}{150}{156}{Solidity}
Questa funzione verifica che il prestito sia ancora finanziato e che il blocco
di scadenza sia stato superato. Poi aggiorna lo stato del prestito, avvisa il
pool e produce un evento. Non è riservata al borrower: i controlli sullo stato
e sul tempo impediscono di dichiarare fallito un prestito non ancora scaduto.

\sourcecode{contracts/FundingPool.sol}{423}{435}{Solidity}
Dentro questo callback \code{msg.sender} è il contratto del prestito.
Il pool verifica che sia autorizzato. \code{isLoanFailed} impedisce di aumentare
più volte il contatore dei fallimenti; \code{loanResolved} impedisce di scalare
più volte quello degli attivi. Alla prima registrazione si aggiustano anche
le percentuali: fallimento $+5$ punti, successo $-5$, nei limiti 1--100\%.

Un rimborso tardivo o una richiesta di compensazione possono anch'essi registrare
il default. Il prestito fallito accetta ancora pagamenti, ma non torna Repaid:
il denaro recuperato dopo non cancella il fatto che Bob fosse in ritardo.

\oral{un prestito risolto è per forza tutto rimborsato?}{No. Risolto significa
che abbiamo registrato l'esito, anche fallimento. Può non esserci più fra gli
attivi e avere ancora capitale non recuperato. Per questo il contatore degli
attivi non basta per sapere quanti ETH Alice può ritirare.}

\chapterpage{15 / LA COMPENSAZIONE}{Recuperare solo ciò che c'è}{La riserva aiuta a coprire una perdita, ma non la garantisce per intero.}
Nella demo Bob non restituisce un prestito di 0,5 ETH. La riserva contiene solo
0,025 ETH, prodotti dal primo prestito andato bene. Prima dell'estratto, la
funzione controlla prestito, scadenza e credito della persona e, se serve,
registra il default.

\sourcecode{contracts/FundingPool.sol}{361}{383}{Solidity}
All'inizio si limita il pagamento al minore tra credito residuo e riserva.
Se il risultato è zero, si emette comunque l'evento e si termina senza ridurre
il credito: Alice potrà riprovare se la riserva viene alimentata di nuovo.

Quando vengono pagati 0,025, diminuiscono \emph{sia} il credito bloccato
\emph{sia} il totale di Alice. Con 50 totali e 0,5 bloccati, il disponibile
è 49,5. Dopo il pagamento resta $49{,}975-0{,}475=49{,}5$.
Il disponibile non sale: gli ETH sono già stati inviati al wallet.

\oral{perché sottrai anche da totalContributed?}{Perché sto pagando fuori dal
pool, non rendendo capitale ritirabile una seconda volta. Se riducessi solo
il bloccato, Alice riceverebbe 0,025 nel wallet e vedrebbe salire di altri
0,025 il saldo da ritirare. Sarebbe un doppio beneficio per lo stesso pagamento.}

\chapterpage{16 / ORACLE SU ETHEREUM}{Una richiesta e una risposta separata}{Il contratto registra il saldo, ma non legge da solo la rete Bitcoin.}
\sourcecode{contracts/BitcoinOracle.sol}{33}{38}{Solidity}
Bob paga almeno \code{minFee} e viene emesso \code{UpdateRequested}.
Qui non viene ancora scritto alcun saldo. L'evento è una richiesta che
il servizio Python deve osservare e servire: pagare non significa ricevere
una risposta sincrona nella stessa chiamata.

\sourcecode{contracts/BitcoinOracle.sol}{45}{48}{Solidity}
Questo è il secondo momento. Il servizio firma con la chiave dell'operatore;
\code{onlyOwner} limita chi può aggiornare il registro. Il numero è in satoshi,
non in BTC decimali e non in wei. Dopo la scrittura un evento segnala la risposta.

\sourcecode{contracts/BitcoinOracle.sol}{54}{56}{Solidity}
Il pool legge il valore memorizzato. Un mapping non valorizzato restituisce
zero: dalla sola lettura zero non posso distinguere indirizzo vuoto da indirizzo
mai aggiornato. La demo aspetta quindi un evento nuovo dopo la propria richiesta.
Anche un evento con saldo zero è una risposta valida, ma non dà capacità di prestito.

\lead{Il costo e il limite di fiducia}
La fee minima è calibrata sul nostro build: 49.426 gas per 0,1 gwei, cioè
0,0000049426 ETH. Gas e prezzo del gas sono due quantità diverse.
L'operatore è fidato: il controllo dell'owner evita scritture non autorizzate,
ma non prova che l'operatore stia dicendo il vero.

\oral{quindi i Bitcoin vengono bloccati a garanzia?}{No. Verifichiamo un saldo
storico, ma non proviamo che Bob controlli quell'indirizzo e non blocchiamo i BTC.
È il modello semplificato del progetto, non una garanzia che il contratto possa
sequestrare in caso di default. Il cambio è fissato a 30 ETH per BTC.}

\chapterpage{17 / IL SALDO BITCOIN}{Contare soltanto gli output non spesi}{Un esempio concreto prima dei dizionari Python.}
Immagina che Alice riceva un output da 2 BTC. Poi lo spenda per pagare Bob
0,4 BTC, ricevendo 1,6 BTC di resto, trascurando la commissione nell'esempio.
Non possiede $2+1{,}6$: il vecchio output da 2 è stato consumato. Ora possiede
solo il nuovo output da 1,6. Un \emph{UTXO} è proprio un output non ancora speso.

\sourcecode{python/oracle/oracle_service.py}{94}{106}{Python}
Per una transazione ordinaria, gli input indicano vecchi output da eliminare.
La coppia \code{spent\_txid, spent\_vout} identifica precisamente quale output
è stato speso. Poi il ciclo sugli output registra i nuovi importi e gli indirizzi
riconoscibili dal parser. \code{None} nella lista delle modifiche significa
``elimina questo output'', non ``il suo valore è zero''.

La coinbase viene trattata separatamente perché crea nuovi output senza spendere
un output precedente ordinario. La ricompensa originale del genesis è esclusa:
non fa parte degli output spendibili. Eventuali donazioni successive allo stesso
indirizzo, invece, possono comparire nel saldo.

\sourcecode{python/oracle/oracle_service.py}{124}{125}{Python}
\code{addr\_utxos} raccoglie gli output ancora presenti per indirizzo.
La funzione somma i loro valori. Se l'indirizzo non è presente, \code{get}
restituisce il dizionario vuoto e la somma vale zero.

\lead{Perché preparare prima le modifiche del blocco?}
La funzione costruisce la lista di rimozioni e inserimenti prima di applicarla.
Se c'è un errore, non vogliamo lasciare mezzo blocco nei registri. Il codice
conserva i vecchi valori dei soli output toccati per poter annullare l'applicazione,
senza copiare ogni volta l'intero insieme UTXO.

\oral{perché non sommi tutte le transazioni ricevute dall'indirizzo?}{Perché
contarei anche denaro che è già stato speso. Per sapere il saldo devo ricostruire
gli output ancora utilizzabili dopo aver applicato entrate e spese in ordine.}

\chapterpage{18 / L'ATTESA DELL'ORACLE}{Che cosa significa ORACLE READY}{Non è una pausa scelta a caso: è la fine della ricostruzione dello snapshot.}
All'avvio l'oracle legge i file Bitcoin locali. Il parser sceglie la catena
collegata al genesis con il lavoro cumulativo maggiore e applica in ordine
i primi \textbf{131.000 blocchi}, dalle altezze 0 a 130.999. Controlla anche hash
e Merkle root, ma non implementa tutte le regole di consenso di un nodo Bitcoin.

\sourcecode{python/oracle/oracle_service.py}{134}{144}{Python}
Il ciclo pretende l'altezza successiva a quella già elaborata. Dopo ogni blocco
incrementa il contatore; alla fine rifiuta una sequenza troppo corta.
Non salta un blocco danneggiato per continuare pubblicando saldi incompleti.
Se la sincronizzazione fallisce, gli indici vengono svuotati e il servizio
non si dichiara pronto.

\sourcecode{python/oracle/oracle_service.py}{195}{201}{Python}
La coda degli eventi viene aperta prima della lettura lunga. Poi si ricostruisce
Bitcoin, si stampa \code{ORACLE READY} e si cominciano a servire le richieste.
Partire dal blocco di deployment evita di dimenticare richieste arrivate
mentre il programma era ancora occupato con Bitcoin.

\lead{Due memorie diverse, da non confondere}
Il saldo Bitcoin viene ricostruito in RAM: a ogni riavvio ricomincia il replay.
SQLite conserva invece eventi Ethereum e transazioni da inviare, non una cache
dei saldi BTC. Una volta pronto, l'oracle risponde usando l'indice in memoria,
senza rileggere tutti i blocchi per ogni richiesta.

\oral{un aggiornamento fresco contiene il saldo BTC di oggi?}{No. È una nuova
risposta ottenuta dallo snapshot storico richiesto. ``Fresco'' nella demo
significa successivo alla nostra richiesta Ethereum, non aggiornato all'ultimo
blocco Bitcoin mondiale. Per questo aspetto READY, non un numero fisso di secondi.}

\chapterpage{19 / IL VOTER AUTOMATICO}{Chi manda davvero il voto YES}{Python esegue una strategia; il contratto applica le regole.}
Il voter legge gli eventi \code{LoanProposed}. Per ogni proposta cerca di
registrare YES con l'account configurato. Non assegna direttamente il prestito
e non salta i controlli di contributo, voto unico e risoluzione.

\sourcecode{python/voter/auto_voter.py}{27}{42}{Python}
Se l'account ha già votato YES, non viene mandata un'altra transazione.
Se aveva votato NO, il programma segnala un errore definitivo: non può cambiare
un voto già registrato. Se la proposta o il servizio sono chiusi, non insiste
come se fosse un problema temporaneo di rete.

Il controllo sul contributo usa il totale, non il disponibile. Se Alice ha
10 ETH interamente bloccati, può ancora votare con peso corrente zero. Se prima
della risoluzione un altro prestito le restituisce capitale, il suo peso potrà
diventare positivo. Questa regola deve essere rispettata sia da Solidity sia
dal programma Python.

\lead{Perché il programma legge prima di scrivere?}
Per capire se la richiesta è ancora da eseguire, se è già stata soddisfatta
o se ormai non può più esserlo. Un'interruzione può avvenire dopo che il voto
è arrivato on-chain ma prima che il programma abbia salvato l'esito locale.
Ricontrollare lo stato evita di trattare quel caso come un nuovo voto.

\oral{dove si vede che la strategia è sempre YES?}{Nella riga 41:
\code{voteOnProposal(proposal\_id, True)}. Il booleano passato è vero.
I controlli prima non scelgono una proposta conveniente: verificano solo che
il voto possa ancora essere registrato correttamente. La maggioranza resta
calcolata dal contratto, usando tutti i contributor.}

\chapterpage{20 / EVENTI E RIAVVII}{Come evitiamo di perdere una richiesta}{Pensala come una lista di lavori da fare e di risposte già preparate.}
Un servizio non deve fidarsi soltanto della memoria del processo: se si chiude,
quella memoria sparisce. SQLite conserva quali eventi sono arrivati, quali
restano da servire e quali transazioni sono già state preparate.

\sourcecode{python/durable_events.py}{52}{62}{Python}
\code{with self.db} raggruppa inserimento degli eventi e aggiornamento del cursore
in una transazione SQLite. Il cursore indica da quale blocco cercare la prossima
volta. Non devo avanzarlo senza aver salvato i lavori: dopo un crash salterei
richieste che non ho mai servito. \code{INSERT OR IGNORE} evita il doppione
di un evento già presente con lo stesso identificativo.

\sourcecode{python/durable_events.py}{173}{177}{Python}
Qui la risposta viene firmata e salvata \emph{prima} dell'invio. Se cade la
connessione dopo la trasmissione, al riavvio posso cercare la stessa transazione
e, se necessario, ritrasmettere gli stessi byte. Non creo alla cieca una risposta
nuova con un altro nonce, che potrebbe pagare due volte la stessa operazione.

\lead{Nonce e hash, senza confonderli}
Il nonce ordina le transazioni di un account. L'hash identifica quella precisa
transazione firmata. Gli account dei servizi devono essere dedicati: se un altro
programma ne usa i nonce, possiamo incontrare conflitti da ispezionare.
La ricevuta dice poi se la transazione è stata inclusa ed è riuscita.

\oral{SQLite garantisce che tutto avvenga esattamente una volta in ogni rete?}{No.
Rende affidabili i retry nel contesto della catena privata prevista. I checkpoint
di hash segnalano reset o riorganizzazioni; non c'è una procedura automatica
generale di recupero dalle reorg. Di fronte a un guasto RPC non cancello la
coda: controllo prima lo stato del nodo e della transazione salvata.}

\chapterpage{21 / SETUP E TRANSAZIONI}{Dal programma Python alla blockchain}{Creare un account, firmare e aspettare una ricevuta sono passaggi distinti.}
\code{setup.sh} prepara gli strumenti e compila; \code{python/setup.py} crea
gli account e distribuisce i contratti. Il setup genera deployer, contributor
e borrower; l'account prefunded fornito dal corso li finanzia soltanto.

\sourcecode{python/setup.py}{57}{67}{Python}
Ogni account riceve 100 ETH. Il prefunded firma questi trasferimenti, ma non
è l'account che distribuisce oracle e pool. Il nuovo deployer crea oracle,
implementazione e proxy, poi lo script verifica owner e configurazione.
I file \code{.env} locali conservano le chiavi generate: non vanno consegnati
né mostrati in una schermata d'esame.

\sourcecode{python/simulate.py}{27}{35}{Python}
\code{build\_transaction} costruisce i dati per chiamare una funzione del contratto.
\code{value} è l'importo da inviare; il nonce usa anche le transazioni pendenti;
il chain ID è 202526, quello del corso. L'account firma localmente. Il nodo
riceve i byte firmati e il client aspetta la ricevuta, controllandone lo stato.

Un hash restituito dal nodo non basta a dire ``operazione riuscita''.
La transazione potrebbe essere ancora pendente o essere inclusa con revert.
L'ABI serve al client per codificare nomi, argomenti ed eventi del contratto;
il bytecode serve invece quando si distribuisce un nuovo contratto.

\oral{perché non chiami tutto con .call()?}{Perché una lettura \code{.call()}
non salva una modifica sulla blockchain. Per depositare, votare o rimborsare
serve una transazione firmata e inclusa in un blocco. Leggere il saldo con
\code{.call()} è diverso dalla \code{call} Solidity che può inviare ETH.}

\chapterpage{22 / PROXY E UPGRADE}{Cambiare la logica senza cambiare indirizzo}{Lo stato del pool rimane nel proxy; i prestiti già creati restano autonomi.}
Immagina che il proxy sia l'indirizzo della cassa e l'implementazione sia il
manuale di istruzioni. Cambiare manuale non crea una nuova cassa né sposta
automaticamente i fondi. Il proxy delega l'esecuzione mantenendo il proprio
storage e saldo: è l'idea della \emph{delegatecall}. Anche \code{msg.sender}
rimane quello della chiamata al proxy: non è una normale chiamata che parte
dal proxy verso un altro contratto.

\sourcecode{contracts/FundingPoolProxy.sol}{9}{12}{Solidity}
Il costruttore riceve l'implementazione e include già la chiamata a
\code{initialize}. Così il proxy nasce inizializzato: non rimane una finestra
in cui un estraneo può chiamare per primo e diventare owner.
\code{ERC1967Proxy} usa uno slot standard per conservare l'indirizzo
dell'implementazione, separandolo dai dati applicativi.

\sourcecode{contracts/FundingPool.sol}{114}{127}{Solidity}
Il costruttore blocca l'inizializzazione dell'istanza di implementazione.
Il proxy inizializza invece il proprio stato tramite la logica: imposta owner,
oracle e collateral iniziale al 50\%. \code{initializer} evita di rifarlo.
L'override alla riga 127 riserva l'upgrade all'owner, sfruttando il meccanismo UUPS.

\lead{Che cosa devo preservare durante un upgrade?}
Il significato e la disposizione dello storage esistente. Se cambio l'ordine
delle variabili senza una migrazione compatibile, la nuova logica potrebbe
leggere vecchi valori come se fossero altri dati. I test V2 controllano alcuni
casi concreti; non rendono sicuro qualunque aggiornamento futuro.

\oral{se aggiorni il pool, cambia anche il codice dei LoanContract?}{No.
Ogni prestito è un contratto autonomo, non un proxy di questo pool. Il suo
bytecode e i termini fissati restano quelli iniziali. Anche per questo una
nuova logica del pool deve continuare a gestire i callback dei vecchi prestiti.}

\chapterpage{23 / REENTRANCY}{Perché il destinatario può richiamarci}{Un trasferimento a un contratto può far eseguire altro codice.}
Il pool normale riduce i registri prima di trasferire ETH. Per mostrare il rischio,
nei test esiste anche una variante volutamente errata: invia prima e aggiorna dopo.
Non viene distribuita dal setup del servizio.

\sourcecode{contracts/test/FundingPoolVulnerable.sol}{14}{23}{Solidity}
Durante la chiamata esterna, il saldo non è ancora stato ridotto. Il contratto
attaccante riceve ETH e richiama \code{withdraw()}; il pool vede ancora il credito
precedente e invia di nuovo. Alla fine vengono riscritti i valori salvati prima
della ricorsione, ormai vecchi.

\sourcecode{contracts/test/ReentrancyAttacker.sol}{30}{37}{Solidity}
La funzione \code{receive()} viene eseguita quando arrivano ETH senza dati.
Se il pool ha ancora denaro, l'attaccante tenta un altro prelievo. La chiamata
a basso livello restituisce un booleano: se il pool normale rifiuta la
ricorsione, l'attaccante può lasciar terminare il prelievo iniziale legittimo.

Nei test l'onesto mette 10 ETH e l'attaccante 1. La variante vulnerabile può
essere svuotata anche dei 10 altrui. Nel pool normale la ricorsione fallisce
e rimangono i fondi dell'onesto.

\oral{bastava spostare una sottrazione dopo la call?}{Non necessariamente.
In Solidity 0.8 una sottrazione sul saldo corrente può andare in underflow
dopo una chiamata ricorsiva e annullare tutto l'attacco. La nostra variante
usa assegnazioni con valori vecchi per mostrare un furto effettivo. Spiego
questo codice preciso, non sostengo che ogni variante apparentemente simile
sia sfruttabile allo stesso modo.}

\chapterpage{24 / TEST E GAS}{Come leggere una prova}{Un test prepara una situazione, esegue un'azione e controlla il risultato.}
\sourcecode{test/trace_compliance.test.ts}{20}{31}{}
\code{deployProtocol()} prepara contratti e account isolati. Tre persone
depositano 2, 3 e 6 ETH; poi viene richiesto un prestito di 4. \code{await}
aspetta il completamento delle operazioni asincrone del test.
Le ultime due righe controllano che il capitale sia inferiore alla richiesta
per arrotondamento e coincida con la somma delle quote calcolate per difetto.
\code{BigInt}, riconoscibile dal suffisso \code{n}, permette conti interi.

\sourcecode{test/gas_profiling.test.ts}{12}{16}{}
Questo helper aspetta la ricevuta e registra \code{gasUsed}, il gas consumato.
Non usa il limite massimo di gas inviato con la transazione.
Il costo in wei si ottiene moltiplicando il gas consumato per il prezzo del gas.
Per esempio 49.426 gas per 0,1 gwei fanno 4.942.600.000.000 wei.

I risultati riportati nel progetto sono 32 test contratti, di cui 4 Solidity
e 28 TypeScript, e 49 test Python offline. Uno scenario Solidity esegue
256 casi fuzz, cioè varia gli input per controllare una proprietà.
Non sono 256 ulteriori test distinti. Il profiling copre 20 scenari.

\lead{Che cosa non dimostra un singolo test}
Il test isolato può minare blocchi immediatamente; la demo Geth no.
Un mock Python non prova che tutto il dataset reale venga letto correttamente.
Il gas con un contributor non è una stima universale per mille contributor.
Il target del compilatore è Berlin, compatibile con il genesis del corso.

\oral{dove trovo la prova di una regola specifica?}{Cerco il test per scenario:
\code{trace\_compliance.test.ts} per i comportamenti di traccia,
\code{requirements\_boundaries.test.ts} per i limiti,
\code{reentrancy.test.ts} per l'attacco e \code{upgrade.test.ts} per il proxy.
Posso spiegare preparazione, azione e controllo, non solo dire ``i test passano''.}

\chapterpage{25 / REGOLE ED ECONOMIA}{Un comportamento scorretto senza bug}{Rispettare il codice non rende ogni incentivo automaticamente sicuro.}
Supponiamo che una persona malevola metta 6 ETH e una persona onesta 4.
La prima controlla la maggioranza e si accorda con il borrower. Può far approvare
un prestito di 10 ETH, purché superi i controlli, e il gruppo può poi non rimborsarlo.

Il gruppo riceve 10, ma 6 erano già capitale proprio: il guadagno netto ai danni
dell'onesto è 4, trascurando gas e altri flussi. Dire che guadagna 10 sarebbe
contare come profitto anche denaro che possedeva prima.

Il dato BTC non impedisce questo comportamento: non richiediamo prova del
possesso dell'indirizzo né blocchiamo i Bitcoin. Inoltre la priorità di recupero
del capitale favorisce chi aveva la quota iniziale più grande. Sono conseguenze
delle regole economiche, non un attacco di reentrancy.

\lead{Perché non cambiare subito il sistema di voto?}
Si potrebbero immaginare altre garanzie, ma il progetto segue la traccia.
Cambiare il peso corrente in uno snapshot di voto, introdurre un vero vincolo
su Bitcoin o redistribuire il rischio in altro modo significherebbe cambiare
il protocollo. All'orale distinguo ciò che abbiamo implementato da possibili
evoluzioni, senza attribuire al progetto protezioni che non possiede.

\sourcecode{contracts/FundingPool.sol}{133}{137}{Solidity}
Questa è la terminazione amministrativa. Solo l'owner può attivarla e non
devono esserci prestiti attivi. \code{terminated} blocca il nuovo business,
non distrugge il contratto. Restano le operazioni previste di prelievo,
incasso e recupero. Un default risolto può comunque lasciare un credito scoperto.

\oral{se terminate vale true, gli ETH spariscono?}{No. È un flag che alcune
funzioni controllano tramite \code{notTerminated}. Non è un selfdestruct e
non rende liquido il capitale non recuperato. Devo guardare quali funzioni
usano il modifier, non immaginare che il flag spenga ogni singola operazione.}

\chapterpage{26 / LA DEMO CONCRETA}{Prima il successo, poi il default}{Collegare i messaggi sul terminale alle funzioni già studiate.}
La demo si avvia con oracle e voter già in esecuzione. Prima aspetta una
risposta BTC nuova, poi finanzia la posizione del contributor. Nel percorso
standard, su deployment nuovo, il deposito è 50 ETH e il primo prestito è 1 ETH
al 5\%. Due rate da 0,525 lo chiudono: gain e riserva ricevono 0,025 ciascuno.
Il gain viene incassato e il collateral globale scende da 50 a 45\%.

\sourcecode{python/simulate.py}{268}{280}{Python}
Il secondo prestito usa metà del capitale richiesto e dura tre blocchi.
\code{max(1, principal // 2)} evita zero per importi minuscoli; \code{//} è
divisione intera Python. \code{wait\_for\_expiry} aspetta un blocco strettamente
maggiore di \code{dueBlock}. La richiesta di compensazione registra il default
e paga il minore fra credito e riserva. Poi la demo ripete la richiesta per
mostrare il comportamento con riserva vuota.

\begin{tabularx}{\linewidth}{@{}Xr@{}}
\toprule Stato finale del percorso standard & Valore\\\midrule
Primo prestito / secondo prestito & Repaid / Defaulted\\
Compensazione pagata / secondo tentativo & 0,025 / 0 ETH\\
Posizione totale del contributor & 49,975 ETH\\
Disponibile / bloccato residuo & 49,5 / 0,475 ETH\\
Riserva / gain da incassare & 0 / 0 ETH\\
Prestiti attivi / falliti & 0 / 1\\
Collateral globale & 50\%\\\bottomrule
\end{tabularx}

Questi numeri descrivono un deployment nuovo con un contributor. Su uno già
usato, altri saldi e riserve possono cambiare i risultati. I wallet pagano anche
il gas, quindi non vanno confrontati con somme che lo ignorano. La prova Geth
salvata in \code{docs/default-demo-verification.json} registra 20 ricevute riuscite
e 16 controlli finali per i sorgenti identificati dagli hash.

\oral{perché il bloccato non torna a zero alla fine?}{Perché il secondo prestito
non è stato rimborsato. Solo 0,025 dei suoi 0,5 ETH sono stati recuperati dalla
riserva. I rimanenti 0,475 sono un credito non recuperato, non ETH ritirabili.}

\chapterpage{27 / PREPARARSI ALLA LETTURA}{Se il professore indica una riga}{Parti dal contesto, poi spiega l'effetto e il motivo.}
Non iniziare traducendo ogni simbolo se non hai ancora detto dove ti trovi.
Puoi usare questo ordine: ``Questa funzione viene chiamata da\ldots;
serve a\ldots; prima controlla\ldots; poi aggiorna\ldots;
questa scelta evita\ldots''. Un piccolo esempio numerico rende la risposta
molto più chiara di un elenco di parole tecniche.

\qa{Ti indica: totalContributed - lockedContributions.}{``Sto calcolando il
capitale libero della persona. Se ha 10 ETH totali e 4 prestati, ne può
usare o ritirare 6. Non sto leggendo il saldo del suo wallet.''}

\qa{Ti indica: proposal.noVotes = currentDisposable - currentYesVotes.}{``Tutto
il capitale disponibile non associato a un YES viene contato come NO.
Comprende anche chi non ha votato. I pesi sono quelli attuali alla risoluzione.''}

\qa{Ti indica: initialLoanContributions e loanContributions.}{``La prima conserva
il contributo iniziale per priorità e interessi; la seconda è il credito
residuo, che diminuisce con rimborsi e compensazioni. Non sono due fondi separati.''}

\qa{Ti indica: !loanResolved[msg.sender].}{``Sto evitando di trattare due volte
la chiusura dello stesso prestito. Il contatore degli attivi non deve scendere
a ogni callback successivo. Qui msg.sender è il contratto del prestito.''}

\qa{Ti indica: if (actualCompensation == 0) return.}{``Il credito è valido, ma
non c'è riserva. Il default è stato registrato prima; esco senza cancellare il
credito. In futuro si può riprovare se arrivano fondi nella riserva.''}

\qa{Ti indica: balances[btcAddress] = balanceInSatoshi.}{``Salvo in Ethereum
la risposta pubblicata dall'operatore. Non ricostruisco Bitcoin in questa riga;
quel lavoro è nel programma Python. Solo l'owner può aggiornare questo dato.''}

\qa{Ti indica: store.save\_transaction prima di send\_raw\_transaction.}{``Salvo la
risposta firmata prima di inviarla. Se il processo si interrompe, posso cercare
e ritentare quella stessa transazione senza crearne una nuova alla cieca.''}

\begin{keybox}{L'ultimo ripasso}
Apri a caso una delle funzioni mostrate. Spiega un caso riuscito e uno rifiutato,
chi è il chiamante e che cosa cambia nei saldi. Se dici ``è per sicurezza'',
aggiungi quale problema preciso evita. Se dici ``lo richiede la traccia'',
distingui la regola dalla nostra scelta di implementazione. Non occorre ricordare
il numero di riga: occorre sapersi orientare nel flusso del prestito.
\end{keybox}
\end{document}
